\section{A logical obstruction and its correct repair}

Replacing the first head
\[
 O(a,H)=(\exists x_a\,\forall x_a H)\to H
\]
by the positive Brouwer head
\[
 B(a,H)=H\to\forall x_a\,\exists x_a H
\]
requires care: the numerical negation is $N(A)=I(A,1)$, and code $1$ is a
membership atom, not a logical falsity constant built into the proof rules.
The old first head supplies explosion. An equivalence argument that first
uses explosion to justify the replacement is circular unless explosion is
separately retained or independently recovered.

\paragraph{Exact obstruction.}
There is a Boolean interpretation of the numerical calculus with every
remaining head, diagonal vacuity, and $B$, in which an instance of $O$ is
false. Interpret $I$ as material implication, every $U(a,\cdot)$ as the
identity, every equality atom as false, $M(0,0)$ as true, and every other
membership atom as false. Interpret code zero as true. This is a recursive
interpretation of the actual tagged numerical formulas; the alias
$I(0,0)=0$ is respected because true implying true is true.

The designated code $1$ is true, so every numerical negation is true.
Consequently every guarded gate and every existential witness head is
true. Membership-congruence heads hold because their equality premises
are false. Commutation, distribution, diagonal vacuity, the restricted
implicational head, and $B$ all hold. Modus ponens and generalization
preserve truth. For a false membership formula $H$, however,
$\exists x_a\forall x_aH$ is true while $H$ is false. Thus $O(a,H)$ fails.
The same interpretation refutes the proposed replacement by
$U(a,H)\to H$ alone. In particular the diagonal clause does not by itself
repair the missing explosion argument.

\paragraph{Repair by an explicit explosion head.}
Let $E(C)=1\to C$. The presentations with $O$ and with both $B,E$ are
uniformly interderivable over the common implicational and quantifier
heads. The old first head gives $E$: under an assumption $1$, use $K$ to
obtain $U(a,N(U(a,C)))\to1$, which is $\exists x_aU(a,C)$; apply $O$ and
discharge the assumption. Both presentations now have explosion and
therefore ordinary double-negation elimination relative to code $1$.

Put $Y=U(a,N(U(a,H)))$. Then $O(a,H)$ is $N(Y)\to H$.
Contraposition followed by double-negation elimination gives
$N(H)\to Y$. Substitute $H=N(C)$ and compose with
$C\to N(N(C))$ to obtain $B(a,C)$.
Conversely, $B(a,N(H))$ gives
\[
 N(H)\to U\bigl(a,N(U(a,N(N(H))))\bigr).
\]
The closed double-negation equivalence $N(N(H))\leftrightarrow H$ lifts
through each displayed quantifier by distribution and generalization of
closed theorems. It yields $N(H)\to Y$. Contraposition and a final
double-negation elimination produce $N(Y)\to H$, as required.
All these derivations are fixed finite proof templates. Bounded scans of
proof codes instantiate them and translate whole proofs; induction on the
input proof length proves their correctness in PA.

Adding a separate explosion head changes the command stack interface by
one entry. Since $E(C)$ is an implication code, it never equals one. Its
presence therefore causes no accidental Shift halt, but every selector,
push/pop word and all-store compiler must use the new head count. A source
or decoder candidate that implements $B$ without this repair remains a
rejected diagnostic candidate for this equivalence argument.

